\documentclass[11pt,letterpaper]{article}
\usepackage{physical_ai_legislation}
\usepackage[backend=biber,style=numeric,sorting=none]{biblatex}
\addbibresource{sf_city_regulations.bib}

\title{%
  \textbf{San Francisco Municipal Code Update}\\[6pt]
  \large Physical AI Oncology Clinical Trial Site\\
  Permit, Zoning, and Operational Requirements\\[6pt]
  \normalsize Proposed amendments to the San Francisco Planning Code,\\
  Building Code, and Health Code for Physical AI trial site authorization.
}
\author{%
  Kevin Kawchak\\
  CEO, ChemicalQDevice\\
  \texttt{kevink@chemicalqdevice.com}
}
\date{March 24, 2026}

\begin{document}
\maketitle
\thispagestyle{fancy}

\begin{center}
\footnotesize
\textit{Unofficial independent draft derived from prior legislation;
no third-party endorsement, sponsorship, affiliation, or authorization
is expressed or implied; and was adapted using Claude Code Opus 4.6.}
\end{center}

\vspace{0.5em}

\begin{center}
\footnotesize
\textit{This draft is independent and is not endorsed, sponsored, or approved
by any trial sponsor, CRO, site, IRB, regulator, or medical society;
and was adapted using Claude Code Opus 4.6.}
\end{center}

\vspace{1em}
\tableofcontents
\newpage

\section{Purpose and Scope}

\subsection{Purpose}

This regulation establishes the municipal permit, zoning, and operational
requirements for Physical AI oncology clinical trial sites within the City and
County of San Francisco, in furtherance of the California Physical AI Oncology
Clinical Trial Authorization and Site Establishment Act (SB 1042). San
Francisco shall be the location of the first Physical AI oncology clinical trial
site in California, situated in a prominent and safe location that provides
accessible care to oncology patients throughout the city and surrounding Bay
Area communities.

\subsection{Scope}

These requirements apply to the permitting, construction, and operation of
new-construction facilities specifically designed as Physical AI oncology
clinical trial sites, including the associated parking facilities required under
this regulation. These requirements supplement but do not replace existing San
Francisco Planning Code, Building Code, and Health Code requirements for medical
facilities.

\subsection{Evidence Base}

The 24-hour on-demand Physical AI oncology clinical trial simulation conducted
on March 23, 2026, provides the primary evidence base for these regulations,
demonstrating that a single 85,000 to 100,000 square-foot facility equipped
with 10 robot types (29 total instances) served 168 unique patients across 15
cancer types with 99.7 percent facility uptime and zero patient
harm~\cite{newtrial2026}. The simulation demonstrated that both specific
clinical details at minute-level resolution and sustained operations across a
full 24-hour cycle were managed correctly by Physical AI
systems~\cite{newtrial2026}.

\section{Zoning Requirements}

\subsection{Permitted Zones}

\begin{enumerate}[label=(\alph*)]
\item Physical AI oncology clinical trial sites shall be a conditionally
  permitted use in the following San Francisco zoning districts:
  \begin{enumerate}[label=(\arabic*)]
  \item Neighborhood Commercial (NC) districts with parcels of 20,000 square
    feet or larger.
  \item Mixed Use-General (MUG) districts.
  \item Production, Distribution, and Repair (PDR) districts adjacent to
    residential or commercial zones with transit access.
  \item Downtown General (DTG) and Downtown Office (C-3) districts.
  \item Medical Use Districts and Health Care Services Districts.
  \end{enumerate}

\item Preference shall be given to locations within one-quarter mile of a Bay
  Area Rapid Transit (BART) station or San Francisco Municipal Railway (Muni)
  stop to maximize patient accessibility consistent with the on-demand 24/7
  scheduling model~\cite{formatcomparison2026}.

\item The site shall be located in a prominent and safe location, defined as a
  parcel with frontage on a street classified as an arterial or collector in
  the San Francisco General Plan, with adequate street lighting, pedestrian
  infrastructure, and proximity to emergency services.
\end{enumerate}

\subsection{Setback and Height Requirements}

\begin{enumerate}[label=(\alph*)]
\item The minimum lot size for a Physical AI oncology clinical trial site shall
  be 40,000 square feet, accommodating the 85,000 to 100,000 square-foot
  building footprint across multiple stories plus dedicated parking
  facilities~\cite{sitespec2026}.

\item Building height shall comply with the applicable district height limit
  but shall not exceed 65 feet unless a variance is obtained through the
  Planning Commission conditional use process.

\item Street-level setbacks shall provide a minimum 15-foot pedestrian zone
  along all public street frontages.

\item The site shall include a dedicated patient drop-off and pick-up zone
  accessible from the public right-of-way, with a covered canopy extending a
  minimum of 20 feet from the building entrance.
\end{enumerate}

\section{Conditional Use Permit Requirements}

\subsection{Application Contents}

An application for a conditional use permit for a Physical AI oncology clinical
trial site shall include, in addition to standard Planning Code requirements,
the following:

\begin{enumerate}[label=(\alph*)]
\item A Physical AI System Deployment Plan identifying all 10 robot types to be
  deployed, the number of instances of each type, and the USL and PSL scores
  for each system, demonstrating compliance with the minimum thresholds
  established by state law~\cite{usl2026,pslframework2026}.

\item A Patient Flow Analysis based on the simulation data, including projected
  daily patient volumes (up to 168 patients per 24 hours), arrival and
  departure patterns, peak utilization periods (the simulation showed a peak
  of 15 arrivals in hour 09 with 72.4 percent utilization and 28 concurrent
  patients on-site), and measures to manage queuing and
  wait times~\cite{newtrial2026}.

\item A Neighborhood Impact Assessment addressing noise, traffic, parking
  demand, emergency vehicle access, and the visual impact of the facility on
  the surrounding neighborhood.

\item A Community Benefits Plan describing workforce development programs,
  local hiring commitments, and patient access programs for San Francisco
  residents.

\item A Transportation Demand Management (TDM) Plan meeting the requirements
  of San Francisco Planning Code Section 169, reflecting the 24/7 operating
  schedule and the reduced staffing model (8 to 10 full-time equivalents
  versus 80 to 120 in traditional trial sites)~\cite{sitespec2026}.
\end{enumerate}

\subsection{Public Notice and Hearing}

\begin{enumerate}[label=(\alph*)]
\item The Planning Department shall provide notice of the conditional use
  permit application to all owners and occupants of property within 300 feet
  of the proposed site.

\item A public hearing before the Planning Commission shall be held not fewer
  than 30 days after notice is provided.

\item The applicant shall hold at least one community meeting in the
  neighborhood where the site is proposed, at least 14 days before the
  Planning Commission hearing.
\end{enumerate}

\section{Building Permit Requirements}

\subsection{Structural and Mechanical Systems}

\begin{enumerate}[label=(\alph*)]
\item Building permit applications for Physical AI oncology clinical trial
  sites shall demonstrate compliance with the California Building Standards
  Code (Title 24) as supplemented by the San Francisco Building Code, with
  the following additional requirements specific to Physical AI operations:

\item Surgical suites shall incorporate HEPA filtration (ISO Class 7 or better),
  positive-pressure air handling, and floor load capacity of at least 150
  pounds per square foot to support surgical robot systems and associated
  equipment~\cite{sitespec2026}.

\item Radiotherapy vaults shall include radiation shielding meeting the
  requirements of the California Radiologic Health Branch, with additional
  structural reinforcement to support the weight of radiotherapy positioning
  robot systems~\cite{sitespec2026}.

\item Server rooms for AI computation infrastructure shall be located on a
  separate HVAC zone with redundant cooling capable of rejecting a minimum
  heat load of 200 kilowatts.

\item The building shall accommodate 800 to 1,200 kilowatts of sustained
  electrical load with redundant utility feeds, automatic transfer switches,
  and uninterruptible power systems covering all critical clinical and
  data systems~\cite{sitespec2026}.
\end{enumerate}

\subsection{Accessibility}

\begin{enumerate}[label=(\alph*)]
\item The site shall exceed the minimum accessibility requirements of the
  Americans with Disabilities Act and California Building Code Chapter 11B
  to ensure that the on-demand patient scheduling model is fully accessible
  to patients with mobility, visual, auditory, or cognitive impairments.

\item All patient circulation routes shall be designed to accommodate hospital
  beds, wheelchairs, and rehabilitation exoskeletons simultaneously.

\item Wayfinding systems shall include visual, auditory, and tactile elements
  appropriate for a facility where patients interact with 10 categories of
  robotic systems~\cite{patientinstructions2026}.
\end{enumerate}

\section{Health Code Requirements}

\subsection{Operating License}

\begin{enumerate}[label=(\alph*)]
\item A Physical AI oncology clinical trial site shall obtain a San Francisco
  Health Care Facility Operating License in addition to the state authorization
  required under SB 1042.

\item The operating license shall specify the 10 categories of Physical AI
  systems authorized for use at the site, the maximum number of concurrent
  patients, and the operating schedule (24 hours per day, 7 days per week).

\item The license shall be renewed annually, contingent upon maintaining a
  cumulative PSL score within the Advanced Site band (60.0 or above) and
  submitting the annual safety and performance
  reports~\cite{pslframework2026}.
\end{enumerate}

\subsection{Inspection Requirements}

\begin{enumerate}[label=(\alph*)]
\item The San Francisco Department of Public Health shall conduct an initial
  pre-operational inspection of the facility before the first patient is
  enrolled.

\item Routine inspections shall occur at least quarterly, with additional
  inspections triggered by adverse event reports, cybersecurity incidents, or
  patient complaints.

\item Inspections shall include verification of Physical AI system safety
  certifications, review of adverse event logs and resolution records, audit
  trail integrity checks, and observation of at least one active procedure.

\item Inspection findings shall be posted on the department's public website
  within 30 days of the inspection date.
\end{enumerate}

\section{Noise and Environmental Standards}

\begin{enumerate}[label=(\alph*)]
\item The site shall comply with the San Francisco Noise Ordinance (Article 29
  of the Police Code) with the following additional provisions for 24-hour
  Physical AI operations:
  \begin{enumerate}[label=(\arabic*)]
  \item HVAC and mechanical equipment noise at the property line shall not
    exceed 45 dBA during nighttime hours (10:00 PM to 7:00 AM).
  \item Loading and delivery activities shall be restricted to daytime hours
    (7:00 AM to 8:00 PM) except for emergency medical supply deliveries.
  \item Robot maintenance and testing activities generating noise above ambient
    levels shall be conducted in acoustically isolated maintenance bays.
  \end{enumerate}

\item Emergency generator testing shall be limited to daytime hours and shall
  not exceed one test per month lasting no more than 30 minutes.

\item The site shall implement a stormwater management plan compliant with San
  Francisco's Green Infrastructure Ordinance, incorporating the parking
  facility's impervious surface area.
\end{enumerate}

\section{Relationship to State Requirements}

\begin{enumerate}[label=(\alph*)]
\item These municipal requirements supplement and do not replace the
  state-level requirements established under SB 1042, AB 2847, and SB 892.

\item Where a municipal requirement exceeds the corresponding state requirement,
  the municipal requirement shall apply. Where a state requirement exceeds the
  corresponding municipal requirement, the state requirement shall apply.

\item The three adapted regulatory frameworks (ICH E6(R3), 21 CFR Part 50, and
  21 CFR Part 312) apply to Physical AI oncology clinical trial sites in San
  Francisco as they apply throughout the state~\cite{iche6r3adapt,cfr50adapt,cfr312adapt}.
\end{enumerate}

\printbibliography[title={References}]

\end{document}
